\chapter{Kết luận và Hướng phát triển}
\label{ch:conclusion}

Chương này tổng hợp hành trình nghiên cứu của luận văn ``Kiến trúc Nhận thức Hai tầng Tự động hóa Phát hiện Đe dọa và Phản hồi Sự cố dùng AI Tác tử'' (SENTINEL). Nó ánh xạ các kết quả đã chứng minh về các câu hỏi nghiên cứu, tự đánh giá giới hạn kiến trúc và thực tiễn, và vạch hướng tương lai để bảo vệ hạ tầng hiện đại bằng AI Tác tử cục bộ nối-nền-ngữ-cảnh.

\section{Tổng hợp Kết quả Nghiên cứu Đạt được}
Kết quả xác nhận giải quyết thành công ba câu hỏi nghiên cứu.

\textbf{RQ1 (kinh tế học của tầng lọc).} Chuỗi lọc hai chặng hấp thụ phần áp đảo lưu lượng trước khi cần tới suy luận nhiều giây. Duy trì baseline per-IP lăn và tính Z-score trong thời gian và bộ nhớ $O(1)$, Tầng 1 tự quyết phần lớn lành tính; \emph{Cổng LightGBM} sau đó tự quyết $83.8\%$ số ca còn leo thang mà không gọi LLM, ở precision $98.82\%$ trên chính phần nó cắt. Mỗi quyết định đường-nhanh hoàn tất dưới một mili-giây (cổng thêm $\approx\!0.35$\,ms) so với một suy luận nhiều giây, tạo ra mức giảm $82.97\%$ độ trễ đầu-cuối trung bình so với pipeline chỉ-LLM. \emph{Cái giá phải trả} được nêu thẳng: cổng đạt recall $0.7551$ trên tập kiểm định, tức phần lưu lượng nó hấp thụ đánh đổi độ phủ lấy thông lượng---đây là đánh đổi trung tâm mà RQ1 đặt ra, không phải một kết quả phụ.

\textbf{RQ2 (phân loại an toàn, độ trễ thấp).} Mô hình \textit{Gemma-2-9B-IT} Q6\_K lượng tử hóa cục bộ được điều khiển bởi FSM LangGraph giải quyết bài toán độ trễ vận hành và lỗ hổng nhận thức. Chỉ các dị biệt đã lọc trước mới chạm tới LLM, và một Semantic Cache Tầng 1.75 trong bộ nhớ (LRU khóa theo băm SHA-256 của mẫu) giải quyết các cảnh báo lặp lại trong $\approx\!1$\,ms, trong khi các dị biệt chưa cache tốn trung bình $4$--$6$\,s mỗi lô, được bổ trợ bởi KV-Caching prompt giúp giảm thêm TTFT cục bộ. Về bảo mật, Đóng gói Dữ liệu Phân định (nonce \texttt{<<<DATA\_HASH...>>>}) cùng với bộ làm sạch đầu ra chặn đứng các kiểu prompt injection qua log (S1--S4) và delimiter smuggling; vì rào chắn đồng thuận là tất định, toàn bộ pipeline kháng lại mọi nỗ lực hạ cấp trong bộ payload thử nghiệm.

\textbf{RQ3 (trạng thái mang lại năng lực gì).} Trạng thái tích lũy sinh ra một năng lực mà tầng phi trạng thái không thể có. \emph{(a)~Liên kết đa ngày nổi lên:} Bộ nhớ Đe dọa bền vững phát hiện $3/3$ chiến dịch APT trải nhiều ngày (recall $1.00$), và---quan trọng hơn---\emph{đối chứng âm} cho specificity $1.00$: bốn IP lành tính cũng hoạt động qua $\ge 2$ ngày phân biệt đều không kích hoạt cờ APT nào. Recall một mình vô nghĩa ở đây (một bộ đếm ngây thơ hô ``APT'' cho mọi IP đa ngày cũng đạt $3/3$); chính đối chứng âm chứng minh hệ \emph{phân biệt} được. \emph{(b)~Quy kết và biện minh:} truy xuất dày FAISS và thưa BM25 hợp nhất bằng RRF ($k=60$) đạt exact-match kỹ thuật $64.0\%$ trên lớp tấn công ứng dụng---trong khi trần phủ của kho tri thức là $80.0\%$---và một LLM-làm-Trọng-tài chéo họ (Meta Llama-3-8B chấm Gemma-2-9B, $n=908$) chấm Faithfulness $3.30$, Answer Relevancy $3.56\pm1.25$, Context Recall $3.27\pm0.95$, Context Precision $2.25\pm0.57$ (tổng $3.1/5$). Precision truy xuất là mắt xích yếu và được nêu ở Mục~\ref{sec:limitations}. Về kiểm chứng thống kê, Mann--Whitney~U xác nhận tách biệt độ trễ hai tầng rất có ý nghĩa ($p<0.001$); McNemar trên tập thiên tấn công cho $p=1.0$, và luận văn \emph{không} diễn giải điều này thành ưu thế phát hiện của Tầng 2---xem thảo luận về giới hạn dụng cụ đo ở Mục~\ref{sec:limitations}.

\section{Đóng góp Khoa học và Thực tiễn}
\textbf{Khoa học.} (1)~Một kiến trúc bảo mật nhận thức hai tầng tách lọc tốc-độ-cao tầng-truyền-tải (Tầng 1, thống kê tất định cộng một cổng giảm-tải học) khỏi suy luận có trạng thái tầng-nhận-thức (Tầng 2), một khuôn mẫu kết hợp toán học tất định với mô hình ngôn ngữ phi tất định. (2)~Một mô hình cách ly mật mã, Đóng gói Dữ liệu Phân định, giải quyết lỗ hổng LLM trước nền-log không tin cậy mà không cần rào chắn phân loại mong manh, tốn tính toán. (3)~Một khung đánh giá phi tham số cho tác tử an ninh cục bộ dùng trọng tài LLM chéo họ và metric RAGAS.

\textbf{Thực tiễn.} (1)~Một bản mẫu mã nguồn mở, đóng gói container bằng Python điều phối qua Docker Compose, tích hợp \texttt{llama.cpp} và \texttt{LangGraph}. (2)~Kiểm chứng phân loại air-gapped cục bộ trên phần cứng phổ thông (i7-14700KF, 32\,GB RAM, một RTX~4060~Ti), với một Cổng LightGBM giảm tải $83.8\%$ ca leo thang ở precision $98.82\%$ để giữ suy luận cục bộ khả thi. (3)~Một dashboard Streamlit nối tự động hóa với hàng đợi Con-người-trong-vòng-lặp, bảo vệ log chống giả mạo bằng chuỗi HMAC.

\section{Các Giới hạn Hiện hữu}
\label{sec:limitations}
Vài giới hạn ràng buộc phạm vi vận hành. \emph{Telemetry:} đánh giá tập trung vào flow mạng và metadata HTTP; các log phi cấu trúc đa dạng (Windows Event, Syslog, CloudTrail) cần thêm cấu hình Drain. \emph{Bão hòa hàng đợi:} một trận quét tốc độ cao hoặc DoS có thể tăng vọt dị biệt và bão hòa hàng đợi leo thang dưới compute đơn-GPU. \emph{Đánh đổi tự chủ:} để tránh tự-từ-chối-dịch-vụ, chặn tường lửa do con người quyết, giới hạn tốc độ giảm thiểu tự chủ. \emph{Đầu độc trực tuyến:} checksum vector ngăn giả mạo offline nhưng không ngăn ghi qua kênh được ủy quyền. \emph{Sức mạnh thống kê:} năng lực zero-day kiểm chứng bằng proxy Welford có kiểm soát (biên $\approx\!4\sigma$) chứ không phải malware thật; các nghiên cứu APT, đối kháng, và chất lượng suy luận dựa trên mẫu khiêm tốn (ba chiến dịch, Wilson $95\%$ CI $[0.44,1.00]$ với đối chứng không-báo-giả; bộ 120 payload; ablation $1{,}250$ mẫu thiên-tấn-công và cân bằng $150/150$, đều trên một mô hình và một host). Bộ đối kháng do tác giả biên soạn, không thích ứng, và khớp mô hình hạ-cấp mà rào chắn đồng thuận vô hiệu, nên là bằng chứng khái niệm, không phải chứng nhận độ bền; đánh giá trước các benchmark injection công khai và tấn công thích ứng là ưu tiên.

\subsection{Giới hạn của Dụng cụ Đo}
Hai giới hạn thuộc về \emph{phép đo}, không thuộc về artifact, và được nêu ở đây vì chúng ràng buộc mạnh cách đọc các con số.

\emph{Thứ nhất, độ phân giải của thước đo nhị phân.} Khi quy mọi phán quyết về một nhãn ``có phát hiện hay không'', các hành động \texttt{ESCALATE} (Tầng 1 chuyển tiếp) và \texttt{AWAIT\_HITL} (tác tử hoãn cho con người) bị gộp cùng ô với \texttt{BLOCK}. Hệ quả là một cấu hình \emph{đẩy quyết định sang tầng sau} được tính điểm ngang một cấu hình \emph{tự quyết đúng}, khiến các cấu hình ablation hội tụ về cùng một giá trị và phép so sánh mất khả năng phân biệt---đây chính là lý do McNemar cho $p=1.0$ ở trên. Vì vậy luận văn \emph{không} rút ra kết luận về đóng góp phát hiện của từng cấu phần từ thước đo nhị phân; các khẳng định về đóng góp cấu phần chỉ dựa trên những phép đo tách biệt (giảm tải cổng ML, quy kết MITRE, đối chứng âm APT, bộ đối kháng). Một thước đo \emph{chấm theo hành động cuối cùng} đối chiếu hành động kỳ vọng---phân biệt tự quyết đúng, tự quyết sai, và hoãn---là điều kiện cần để định lượng đóng góp từng cấu phần, và được nêu như công việc trước mắt.

\emph{Thứ hai, tầng LLM không phải một bộ phân loại.} Khi cố tình bỏ qua cổng ML và đẩy thẳng một tập $95\%$ lành tính vào tác tử ($n=800$), tầng LLM đạt recall $1.00$ nhưng specificity $0.00$, với accuracy $0.0475$ \emph{đúng bằng} mốc đa số---nghĩa là ở chế độ đó nó gắn cờ gần như mọi thứ. Con số này không mô tả đường vận hành thật (nơi cổng ML lọc trước), nhưng nó xác lập một ranh giới trung thực: tầng nhận thức đóng vai một \emph{lưới an toàn thiên về độ phủ} kèm khả năng biện minh, chứ không phải một bộ phân biệt lành/độc. Mọi phát biểu trong luận văn về tầng này được giữ trong ranh giới đó.

\section{Hướng Phát triển Tương lai}
Vài hướng tiếp nối. Một phân rã \emph{đa-tác-tử} sẽ tách tác tử nguyên khối thành các tác tử con hợp tác (trích IOC, liên kết playbook, biên dịch luật, báo cáo pháp y). \emph{Học liên kết} sẽ chia sẻ trọng số hoặc tín hiệu uy tín huấn luyện cục bộ giữa các nút mà không trao đổi log thô. \emph{Học tăng cường} sẽ cập nhật tham số sinh luật từ phê duyệt analyst trong hàng đợi HITL. \emph{Tự-củng-cố đối kháng} sẽ triển khai một tác tử đối thủ trong mạng để dò rào chắn tự động. \emph{Quy kết dựa-TTP} sẽ đặt một module sinh giả thuyết trên các định danh ATT\&CK có cấu trúc, so chuỗi kỹ thuật per-source với hồ sơ MITRE ATT\&CK Group bằng độ tương đồng tập có trọng số---trình bày nghiêm ngặt như một giả thuyết xếp hạng để analyst xem, không bao giờ là phán quyết dứt khoát, vì kỹ thuật được chia sẻ giữa các nhóm và quy kết thực cần tình báo vượt xa nghiệp vụ.

Hai năng lực nền tảng cần \emph{trưởng thành hướng-sản-xuất}. Bộ lọc zero-day Tầng 1 hiện là Z-score Welford đơn-biến per-đặc-trưng (biên $\approx\!4\sigma$); sản xuất sẽ mở rộng nó tới mô hình đa biến hoặc học nhẹ kiểm chứng trước malware thật. Liên kết APT đa ngày hiện khóa theo IP nguồn (bị đánh bại bởi xoay IP hay fan-out botnet); khóa lại theo kỹ thuật và thực thể ổn định sẽ giữ chiến dịch liên kết được qua hạ tầng đổi thay. Về phương pháp, một tái đánh giá APT trung thực sẽ khóa việc ghi Bộ nhớ Đe dọa theo phán quyết hai tầng \emph{trực tuyến} (không theo nhãn DAPT2020) và phân vùng lưu lượng theo thời gian bắt thật, đệm mỗi ngày bằng nền benign quy-mô-sản-xuất---cho recall đầu-cuối thấp hơn nhưng trung thực và bộc lộ báo-giả-APT, qua đó thúc đẩy khóa-lại theo kỹ thuật/thực thể ở trên.

\subsection{Lộ trình Mở rộng Ngang}
Bản mẫu đơn-host thừa hưởng nút thắt đơn-nút. Một tiến hóa sản xuất theo playbook hệ-phân-tán mà không đổi hợp đồng hai tầng: (i)~\emph{phân mảnh bộ lọc Tầng 1 có trạng thái} bằng băm nhất quán trên khóa IP-nguồn, giữ nguyên baseline per-IP và chuỗi APT trong khi năng lực tăng tuyến tính; (ii)~thay broker Redis đơn bằng một \emph{commit log phân mảnh, sao chép} (Redis Cluster hoặc Kafka) có khóa phân mảnh khớp khóa shard Tầng 1; (iii)~thêm \emph{tiêu thụ idempotent} khóa theo định danh sự kiện tất định để một tái-giao sau sự cố không thực thi chặn hai lần; (iv)~mở rộng Tầng 2 thành một \emph{pool đàn hồi} các worker suy luận phi trạng thái autoscale theo độ sâu hàng đợi, trực tiếp giảm bão hòa trong khi Tầng 1 bảo vệ độc lập; và (v)~di trú các kho SQLite sang một \emph{tầng lưu trữ sao chép} (kho key-value quorum cho uy tín/APT, nhất quán mạnh cho audit HMAC và kho luật). Một dashboard phi trạng thái sau load balancer và một rate limiter phía thu nạp hoàn thiện bức tranh. Cho quy mô doanh nghiệp, năm tối ưu cụ thể:
\begin{itemize}
    \item \textbf{Tách microservice:} trích Tầng 2 thành dịch vụ độc lập giao tiếp qua message broker, thoát GIL Python và mở rộng qua nhiều máy chủ.
    \item \textbf{Tầng 1 biên dịch (Go/Rust):} viết lại \texttt{RuleEngine} bằng ngôn ngữ hệ thống để loại chi phí JSON và đạt hàng triệu sự kiện/giây.
    \item \textbf{Lưu trữ Kafka:} thay Redis Streams bằng Kafka ghi đĩa để đệm burst volumetric không rủi ro OOM.
    \item \textbf{Batching vLLM / TGI:} thay \texttt{llama.cpp} bằng continuous batching dựa PagedAttention để phán xử hàng chục ca leo thang mỗi chu kỳ GPU.
    \item \textbf{Độ ổn định SQLite \& Redis:} bắt buộc \texttt{synchronous=NORMAL} (tuyệt đối tránh chế độ \texttt{WAL}) và dùng \texttt{lag} của consumer group để backpressure, ngăn vòng lặp crash cross-UID trong Docker.
    \item \textbf{Ghi APT theo phán quyết:} khi sản xuất, ghi mọi IP Z-cao hoặc vi phạm luật vào Bộ nhớ Đe dọa kèm mốc thời gian, tự gắn cờ APT nổi lên khi vi phạm trải 2--3 ngày.
\end{itemize}

\vspace{0.6cm}
\begin{flushright}
    \textit{Khép lại một nỗ lực khiêm tốn, nhưng mở ra những khát vọng sâu xa trong hành trình bảo vệ không gian số bằng sức mạnh của AI Tác tử.}
\end{flushright}
